\subsection{Comparison by data property}\label{sec:comparison}
\owner{All}

% Rubric: "compare the performance", "strengths and weaknesses", "insightful comparisons".
% Every number must come from a results CSV in the repo. Put the path in a % comment.

\todo{All}{One results table per property (rows: dataset $\times$ method; columns: \#clusters,
noise, S, DB, ARI, NMI, time, fake-data gap). Bold the best per dataset.}

% [Wei Yew]: Table tab:cmp_d2d8 uses the shared columns; merge its rows into the per-property tables
% once the other datasets are in the same format.
% Sources: weiyew/sc4020-earthquake/results/method_comparison.csv; weiyew/sc4020-hdb/results/comparison.csv.
\begin{table}[h]\centering\small
\caption{D2 and D8 at the settings chosen in Section~\ref{sec:protocol}. S: silhouette; DB:
Davies--Bouldin; null S: silhouette of the same setting on fake data (D2: shuffled coordinates for K-means,
uniform points for density methods; D8: shuffled columns). ``none'': no clusters formed on the fake data.
D2 has no labels.}
\label{tab:cmp_d2d8}
\resizebox{\textwidth}{!}{\begin{tabular}{llrrrrrrrrr}
\toprule
\textbf{Data} & \textbf{Method} & \textbf{Clusters} & \textbf{Noise} & \textbf{Largest} & \textbf{S} & \textbf{DB}
& \textbf{ARI} & \textbf{NMI} & \textbf{Time (s)} & \textbf{Null S} \\
\midrule
D2 & K-Means ($K=6$)      & 6  & 0\%    & 44.3\% & 0.525 & 0.800 & -- & -- & 0.004 & 0.389 \\
D2 & K-Means++ ($K=6$)    & 6  & 0\%    & 31.0\% & 0.528 & 0.738 & -- & -- & 0.005 & 0.391 \\
D2 & DBSCAN (101.8\,km, 20) & 41 & 40.4\% & 11.8\% & 0.703 & 0.351 & -- & -- & 0.051 & none \\
D2 & HDBSCAN (50)         & 28 & 30.7\% & 7.5\%  & 0.757 & 0.329 & -- & -- & 1.081 & 0.320 \\
\midrule
D8 & K-Means ($K=4$)      & 4  & 0\%    & 36.2\% & 0.317 & 1.116 & 0.220 & 0.299 & 0.011 & 0.195 \\
D8 & K-Means++ ($K=4$)    & 4  & 0\%    & 35.1\% & 0.317 & 1.116 & 0.213 & 0.293 & 0.010 & 0.218 \\
D8 & DBSCAN (0.567, 10)   & 3  & 1.4\%  & 97.1\% & 0.239 & 0.732 & 0.030 & 0.091 & 0.428 & 0.330 \\
D8 & HDBSCAN (200)        & 2  & 28.5\% & 68.6\% & 0.160 & 1.133 & 0.072 & 0.101 & 2.022 & 0.593 \\
\bottomrule
\end{tabular}}
\end{table}

\subsubsection{Curved and elongated clusters}\label{sec:shape}
\owner{Lam (D1), Wei Yew (D2)}
\todo{Lam}{D1: DBSCAN ARI 0.982, Spectral 1.000, K-Means++ 0.481 (\nolinkurl{lam/results/make_moons_results.csv}).}

% Source: weiyew/sc4020-earthquake/notebooks/ablation.ipynb, section 5 (K-means belt case, printed
% centre distances); figure from the same cell.
K-means assumes compact clusters, and D2 breaks that assumption at planetary scale. Earthquakes follow
plate boundaries, so a cluster is a long,
narrow belt rather than a ball around a centre. K-Means++ with $K=6$ must cover the whole globe with six
roughly round regions, so each region takes in several unrelated belts
(Figure~\ref{fig:kmeans_belts}). One cluster of 905 earthquakes joins Turkey, the Reykjanes Ridge and
Iceland, Afghanistan and Tajikistan, Somalia and Eritrea, and the Mid-Indian Ridge. HDBSCAN finds each of
these as its own cluster, and their centres are 2{,}932 to 11{,}449\,km apart. K-means also has no noise
label, so every isolated earthquake in the oceans is assigned to some region. Its silhouette is still
0.528, well above the 0.391 of shuffled coordinates: the score confirms that the data are not random, not
that the clusters are meaningful.

\begin{figure}[h]\centering
\includegraphics[width=0.9\textwidth]{eq_case_kmeans_merges_belts}
\caption{D2, failure of K-Means++ ($K=6$). All points shown in colour or dark grey belong to one K-means
cluster; colours are the HDBSCAN clusters inside it, crosses their centres. Cause: cluster shape. Belts
thousands of kilometres apart are joined because K-means can only draw compact regions.}
\label{fig:kmeans_belts}
\end{figure}

\subsubsection{Uneven density}\label{sec:density}
\owner{Wei Yew (D2), Ky and Hoang Anh (D3)}
% Sources: weiyew/sc4020-earthquake/results/case_density_regions.csv, method_comparison.csv;
% k-distance spread from weiyew/sc4020-earthquake/notebooks/eda.ipynb, section 12.
On D2 the distance to the nearest earthquake runs from 1.2\,km to 100.5\,km (5th to 95th percentile), and
the k-distance curve rises smoothly with no knee. Any single $\varepsilon$ is therefore too small for some
belts and too large for others. At the chosen $\varepsilon=101.8$\,km, DBSCAN clusters 90\% of the
earthquakes in the dense south-west Pacific but only 20\% in the sparser Americas
(Figure~\ref{fig:density}). HDBSCAN, which picks a density level per cluster, clusters 62\% of the
Americas and keeps the Mexico to Panama belt and the Chile and Argentina belt as long clusters. The price is paid in the dense
region, where it clusters 78\% instead of 90\%. Overall HDBSCAN keeps more earthquakes (69.3\% against
59.6\%) with a higher silhouette (0.757 against 0.703, Table~\ref{tab:cmp_d2d8}).

\begin{figure}[h]\centering
\includegraphics[width=0.85\textwidth]{eq_case_density_dbscan_vs_hdbscan}
\caption{D2, DBSCAN ($\varepsilon=101.8$\,km) against HDBSCAN in a sparse region (top, the Americas)
and a dense one (bottom, the south-west Pacific); grey points are noise. Success of HDBSCAN on the sparse
belts, at the cost of trimming the dense ones. Cause: density varies along and between belts.}
\label{fig:density}
\end{figure}
\todo{Ky + Hoang Anh}{D3: Midtown vs outer boroughs; knee $\varepsilon$ vs silhouette-optimal $\varepsilon$.}

\subsubsection{High dimensions}\label{sec:dims}
\owner{Ky (D4), Lam (D5)}
\todo{Ky}{D4: crossover between 3 and 4 PCA dimensions (Ablation D).}
\todo{Lam}{D5: DBSCAN ARI 0.033; HDBSCAN labels 71\% of points noise.}

\subsubsection{No gaps between groups}\label{sec:nogaps}
\owner{Wei Yew (D8)}
% Sources: weiyew/sc4020-hdb/results/comparison.csv, kmeans_k_sweep.csv, screening.csv.
D8 is one continuous mass: flats come in every combination of size, lease, storey and price, with no empty
space between groups. Table~\ref{tab:cmp_d2d8} shows what each method does with it. K-means still finds a
weak partition, four market segments (newer mid-size, older small, large older, newer high-floor), with a
silhouette of 0.317 against 0.218 on shuffled columns. That ratio of 1.45 is modest, but it is the only
result on D8 that beats its fake data. Density methods have no sparse region to cut along. The DBSCAN
setting chosen by the rule puts 97.1\% of sales in one cluster, and HDBSCAN finds one core holding 68.6\%
of sales plus a compact block of new 3-room flats, and calls 28.5\% noise. Both score below the same
setting on shuffled columns (0.239 against 0.330, and 0.160 against 0.593). On shuffled columns HDBSCAN calls
94.0\% of sales noise and keeps two small, well-separated clusters, which the silhouette rewards. On data
with no gaps, a density method has nothing to find, and the selection rule
cannot tell.

\begin{figure}[h]\centering
\includegraphics[width=0.95\textwidth]{hdb_comparison}
% Three of the six panels of weiyew/sc4020-hdb/figures/comparison.png (recorded data), cropped side by side.
\caption{D8 at the chosen settings on recorded data, drawn on the first two principal components
(6{,}000 sales; grey: noise). K-means cuts the mass into four segments; DBSCAN and HDBSCAN return one dominant cluster. Cause:
no gaps between groups.}
\label{fig:hdb_cmp}
\end{figure}

\subsubsection{Silhouette versus ground truth}\label{sec:silhouette}
\owner{All}
\todo{All}{One paragraph, one line per dataset:
D4 best silhouette (HDBSCAN) has the worst ARI, rank correlation $-0.80$ (Ky);
D6 DBSCAN silhouette 0.42 vs K-Means++ 0.28, ARI 0.31 vs 0.90 (Lam);
D7 K-Means top silhouette 0.343, GMM best ARI 0.774 vs 0.654 (Hoang Anh);
D8 unscaled K-means silhouette 0.540 from price bands, against 0.317 standardised (Wei Yew, written in
Ablation~C); D2 K-means silhouette flat at 0.516 to 0.532 for every $K$ from 4 to 12, so it cannot choose
$K$ (Wei Yew, \nolinkurl{weiyew/sc4020-earthquake/results/kmeans_k_selection.csv});
D5 HDBSCAN top silhouette with 71\% noise (Lam).}

\begin{figure}[h]\centering
\includegraphics[width=0.72\textwidth]{f8_metric_disagreement}
\caption{\todo{Ky}{existing figure; extend to D5--D8 or make one panel per dataset.}}
\label{fig:silhouette}
\end{figure}
